\textbf{Phần 1. Trắc nghiệm lựa chọn}
\setcounterpageref{ex}{0}

\begin{ex} %[3H3V50]
Một khu đất hình chữ nhật có chiều dài $45\text{ m}$, chiều rộng $25\text{ m}$. Người ta xây hàng rào xung quanh khu đất và để một cổng vào rộng $4\text{ m}$. Chiều dài hàng rào là
\choice
{$140\text{ m}$}
{\True $136\text{ m}$}
{$130\text{ m}$}
{$144\text{ m}$}
\loigiai{
Chu vi khu đất: $(45 + 25) \times 2 = 140\text{ m}$. Chiều dài hàng rào: $140 - 4 = 136\text{ (m)}$.
\textbf{Đáp án của câu hỏi là B}
}
\end{ex}

\begin{ex} %[3H3V50]
Một hình vuông có chu vi bằng chu vi hình chữ nhật có chiều dài $28\text{ cm}$ và chiều rộng $12\text{ cm}$. Độ dài cạnh của hình vuông đó là
\choice
{$40\text{ cm}$}
{$16\text{ cm}$}
{$24\text{ cm}$}
{\True $20\text{ cm}$}
\loigiai{
Chu vi hình chữ nhật: $(28 + 12) \times 2 = 80\text{ cm}$. Cạnh hình vuông: $80 : 4 = 20\text{ (cm)}$.
\textbf{Đáp án của câu hỏi là D}
}
\end{ex}

\begin{ex} %[3H3V50]
Một hình tam giác có chu vi $72\text{ cm}$. Biết độ dài ba cạnh của hình tam giác bằng nhau. Độ dài mỗi cạnh của hình tam giác đó là
\choice
{\True $24\text{ cm}$}
{$18\text{ cm}$}
{$36\text{ cm}$}
{$20\text{ cm}$}
\loigiai{
Độ dài mỗi cạnh: $72 : 3 = 24\text{ (cm)}$.
\textbf{Đáp án của câu hỏi là A}
}
\end{ex}

\begin{ex} %[3H3V50]
Một hình chữ nhật có nửa chu vi là $36\text{ cm}$, chiều dài hơn chiều rộng $8\text{ cm}$. Chiều rộng của hình chữ nhật đó là
\choice
{$22\text{ cm}$}
{$18\text{ cm}$}
{\True $14\text{ cm}$}
{$16\text{ cm}$}
\loigiai{
Chiều rộng: $(36 - 8) : 2 = 14\text{ (cm)}$.
\textbf{Đáp án của câu hỏi là C}
}
\end{ex}

\begin{ex} %[3H3V50]
Một mảnh vườn hình vuông có cạnh dài $18\text{ m}$. Người ta rào xung quanh bằng 3 tầng dây thép gai. Tổng chiều dài dây thép gai cần dùng là
\choice
{$72\text{ m}$}
{\True $216\text{ m}$}
{$144\text{ m}$}
{$200\text{ m}$}
\loigiai{
Chu vi mảnh vườn: $18 \times 4 = 72\text{ m}$. Ba tầng dây thép gai: $72 \times 3 = 216\text{ (m)}$.
\textbf{Đáp án của câu hỏi là B}
}
\end{ex}

\begin{ex} %[3H3V50]
Người ta ghép hai hình vuông nhỏ bằng nhau cạnh $6\text{ cm}$ thành một hình chữ nhật. Chu vi của hình chữ nhật nhận được là
\choice
{$48\text{ cm}$}
{$30\text{ cm}$}
{$42\text{ cm}$}
{\True $36\text{ cm}$}
\loigiai{
Hình chữ nhật có chiều dài $12\text{ cm}$, chiều rộng $6\text{ cm}$. Chu vi: $(12 + 6) \times 2 = 36\text{ (cm)}$.
\textbf{Đáp án của câu hỏi là D}
}
\end{ex}

\begin{ex} %[3H3V50]
Một hình tứ giác có độ dài các cạnh lần lượt là $12\text{ cm}, 15\text{ cm}, 18\text{ cm}, 25\text{ cm}$. Chu vi của hình tứ giác đó là
\choice
{\True $70\text{ cm}$}
{$60\text{ cm}$}
{$65\text{ cm}$}
{$75\text{ cm}$}
\loigiai{
Chu vi tứ giác: $12 + 15 + 18 + 25 = 70\text{ (cm)}$.
\textbf{Đáp án của câu hỏi là A}
}
\end{ex}

\begin{ex} %[3H3V50]
Nếu cạnh của một hình vuông tăng thêm $5\text{ cm}$ thì chu vi của hình vuông đó tăng thêm bao nhiêu xăng-ti-mét?
\choice
{$5\text{ cm}$}
{$10\text{ cm}$}
{\True $20\text{ cm}$}
{$25\text{ cm}$}
\loigiai{
Chu vi bằng 4 lần cạnh. Khi cạnh tăng $5\text{ cm}$ thì chu vi tăng: $5 \times 4 = 20\text{ (cm)}$.
\textbf{Đáp án của câu hỏi là C}
}
\end{ex}

\begin{ex} %[3H3V50]
Một hình chữ nhật có chu vi $60\text{ cm}$. Nếu tăng chiều rộng thêm $4\text{ cm}$ và giảm chiều dài đi $4\text{ cm}$ thì được một hình vuông. Cạnh của hình vuông đó là
\choice
{$19\text{ cm}$}
{\True $15\text{ cm}$}
{$11\text{ cm}$}
{$16\text{ cm}$}
\loigiai{
Khi tăng chiều rộng $4\text{ cm}$ và giảm chiều dài $4\text{ cm}$ thì chu vi không đổi (vẫn bằng $60\text{ cm}$). Cạnh hình vuông: $60 : 4 = 15\text{ (cm)}$.
\textbf{Đáp án của câu hỏi là B}
}
\end{ex}

\begin{ex} %[3H3V50]
Một sợi dây thép dài $1\text{ m } 20\text{ cm}$ được uốn thành một hình tam giác có 3 cạnh bằng nhau. Mỗi cạnh của hình tam giác đó dài là
\choice
{$30\text{ cm}$}
{$60\text{ cm}$}
{$45\text{ cm}$}
{\True $40\text{ cm}$}
\loigiai{
Đổi $1\text{ m } 20\text{ cm} = 120\text{ cm}$. Độ dài mỗi cạnh: $120 : 3 = 40\text{ (cm)}$.
\textbf{Đáp án của câu hỏi là D}
}
\end{ex}

\begin{ex} %[3H3V50]
Một hình chữ nhật có chiều dài gấp 3 lần chiều rộng và chu vi bằng $64\text{ cm}$. Chiều rộng của hình chữ nhật đó là
\choice
{\True $8\text{ cm}$}
{$16\text{ cm}$}
{$24\text{ cm}$}
{$6\text{ cm}$}
\loigiai{
Nửa chu vi: $64 : 2 = 32\text{ cm}$. Chiều dài chiếm 3 phần, rộng 1 phần: $3 + 1 = 4$ phần. Chiều rộng là: $32 : 4 = 8\text{ (cm)}$.
\textbf{Đáp án của câu hỏi là A}
}
\end{ex}

\begin{ex} %[3H3V50]
Người ta uốn một đoạn dây đồng vừa đủ thành một hình chữ nhật có chiều dài $16\text{ cm}$ và chiều rộng $10\text{ cm}$. Nếu dùng đoạn dây đó uốn thành hình vuông thì cạnh hình vuông là
\choice
{$14\text{ cm}$}
{$12\text{ cm}$}
{\True $13\text{ cm}$}
{$15\text{ cm}$}
\loigiai{
Chu vi đoạn dây: $(16 + 10) \times 2 = 52\text{ cm}$. Cạnh hình vuông: $52 : 4 = 13\text{ (cm)}$.
\textbf{Đáp án của câu hỏi là C}
}
\end{ex}

\textbf{Phần 2. Trắc nghiệm đúng sai}
\setcounterpageref{ex}{0}

\begin{ex} %[3H3V50]
Xét tính đúng sai của các công thức tính chu vi:
\choiceTF
{\True Chu vi hình chữ nhật bằng (chiều dài + chiều rộng) $\times 2$ (cùng đơn vị đo).}
{Chu vi hình vuông bằng cạnh nhân với chính nó.}
{\True Chu vi hình tam giác bằng tổng độ dài 3 cạnh của nó.}
{Hai hình chữ nhật có chu vi bằng nhau thì chắc chắn có chiều dài và chiều rộng bằng nhau.}
\loigiai{
\begin{itemize}
    \item Ý a) Đúng: Công thức chuẩn.
    \item Ý b) Sai: Sai vì đó là công thức diện tích, chu vi bằng cạnh $\times 4$.
    \item Ý c) Đúng: Định nghĩa chu vi.
    \item Ý d) Sai: Sai vì có thể khác kích thước nhưng cùng tổng chiều dài và rộng.
\end{itemize}
\begin{itemize}
    \item a) \textbf{Đúng}.
    \item b) \textbf{Sai}.
    \item c) \textbf{Đúng}.
    \item d) \textbf{Sai}.
\end{itemize}
}
\end{ex}

\begin{ex} %[3H3V50]
Một sân trường hình chữ nhật có chiều dài $60\text{ m}$, chiều rộng $40\text{ m}$:
\choiceTF
{\True Nửa chu vi sân trường là $100\text{ m}$.}
{\True Chu vi sân trường là $200\text{ m}$.}
{Bạn An chạy một vòng quanh mép sân trường hết một quãng đường dài $100\text{ m}$.}
{\True Nếu giảm chiều dài đi $20\text{ m}$ thì sân trường trở thành hình vuông.}
\loigiai{
\begin{itemize}
    \item Ý a) Đúng: $60 + 40 = 100\text{ m}$.
    \item Ý b) Đúng: $100 \times 2 = 200\text{ m}$.
    \item Ý c) Sai: Sai vì một vòng quanh sân là 1 chu vi = $200\text{ m}$.
    \item Ý d) Đúng: $60 - 20 = 40\text{ m}$ bằng chiều rộng.
\end{itemize}
\begin{itemize}
    \item a) \textbf{Đúng}.
    \item b) \textbf{Đúng}.
    \item c) \textbf{Sai}.
    \item d) \textbf{Đúng}.
\end{itemize}
}
\end{ex}

\textbf{Phần 3. Trả lời ngắn}
\setcounterpageref{ex}{0}

\begin{ex} %[3H3V50]
Tính chu vi hình vuông có cạnh dài $24\text{ cm}$.
\shortans{96}
\loigiai{
\textbf{Đáp số:} $96$\\
$24 \times 4 = 96\text{ (cm)}$.
}
\end{ex}

\begin{ex} %[3H3V50]
Một hình chữ nhật có chu vi $80\text{ cm}$, chiều dài là $26\text{ cm}$. Tính chiều rộng hình chữ nhật theo xăng-ti-mét.
\shortans{14}
\loigiai{
\textbf{Đáp số:} $14$\\
Nửa chu vi: $80 : 2 = 40\text{ cm}$. Chiều rộng: $40 - 26 = 14\text{ (cm)}$.
}
\end{ex}

\begin{ex} %[3H3V50]
Một mảnh vườn hình chữ nhật có chiều rộng $15\text{ m}$, chiều dài gấp đôi chiều rộng. Tính chu vi mảnh vườn đó theo mét.
\shortans{90}
\loigiai{
\textbf{Đáp số:} $90$\\
Chiều dài: $30\text{ m}$. Chu vi: $(30 + 15) \times 2 = 90\text{ (m)}$.
}
\end{ex}

\begin{ex} %[3H3V50]
Cần bao nhiêu mét lưới rào quanh một khu vườn hình vuông có cạnh $35\text{ m}$ nếu chừa lối đi rộng $3\text{ m}$?
\shortans{137}
\loigiai{
\textbf{Đáp số:} $137$\\
Chu vi: $35 \times 4 = 140\text{ m}$. Chiều dài lưới: $140 - 3 = 137\text{ (m)}$.
}
\end{ex}

\begin{ex} %[3H3V50]
Một tam giác có độ dài các cạnh là $18\text{ cm}, 22\text{ cm}$ và $25\text{ cm}$. Tính chu vi hình tam giác đó theo xăng-ti-mét.
\shortans{65}
\loigiai{
\textbf{Đáp số:} $65$\\
$18 + 22 + 25 = 65\text{ (cm)}$.
}
\end{ex}

\begin{ex} %[3H3V50]
Người ta ghép 4 hình vuông cạnh $5\text{ cm}$ thành một hình vuông lớn hơn. Tính chu vi của hình vuông lớn đó theo xăng-ti-mét.
\shortans{40}
\loigiai{
\textbf{Đáp số:} $40$\\
Hình vuông lớn có cạnh: $5 \times 2 = 10\text{ cm}$. Chu vi: $10 \times 4 = 40\text{ (cm)}$.
}
\end{ex}
